
\documentclass{article}
\usepackage{colortbl}
\usepackage{makecell}
\usepackage{multirow}
\usepackage{supertabular}

\begin{document}

\newcounter{utterance}

\centering \large Interaction Transcript for game `cladder', experiment `full\_v1.5\_default', episode 6252 with Qwen3.5{-}4B.
\vspace{24pt}

{ \footnotesize  \setcounter{utterance}{1}
\setlength{\tabcolsep}{0pt}
\begin{supertabular}{c@{$\;$}|p{.15\linewidth}@{}p{.15\linewidth}p{.15\linewidth}p{.15\linewidth}p{.15\linewidth}p{.15\linewidth}}
    \# & \multicolumn{2}{c}{Player} && \multicolumn{2}{c}{Game Master} \\
    \hline

    \theutterance \stepcounter{utterance}  
    & & & \multicolumn{4}{p{0.6\linewidth}}{
        \cellcolor[rgb]{0.9,0.9,0.9}{
            \makecell[{{p{\linewidth}}}]{
                \texttt{\tiny{[P1$\langle$GM]}}
                \texttt{You are an expert in causal inference. The following question is not a typical commonsense query, but rather a meticulously designed question created by a professor specializing in causal inference, intended to assess the students' mastery of the course content.} \\
\texttt{Imagine a self{-}contained, hypothetical world with only the following conditions, and without any unmentioned factors or causal relationships: Unobserved confounders has a direct effect on maternal smoking status and brown eyes. Cigarette tax has a direct effect on maternal smoking status. Maternal smoking status has a direct effect on brown eyes. Unobserved confounders is unobserved. We know that confounder active and high cigarette tax causes smoking mother. confounder active and smoking mother causes brown eyes. We observed the cigarette tax is high and confounder active. Would the infant has blue eyes if smoking mother instead of nonsmoking mother?} \\
\texttt{Start your answer with 'Yes' or 'No', followed by additional reasoning or evidence to support your explanation.} \\
            }
        }
    }
     \\ \\

    \theutterance \stepcounter{utterance}  
    & & \multicolumn{4}{p{0.6\linewidth}}{
        \cellcolor[rgb]{0.95,0.95,0.95}{
            \makecell[{{p{\linewidth}}}]{
                \texttt{\tiny{[P1$\leftrightarrow$P]}}
                \texttt{{[}DRAFT{]}} \\
\texttt{No} \\
\\ 
\texttt{**Reasoning:**} \\
\\ 
\texttt{In this hypothetical world, the causal pathway from maternal smoking status to the infant's eye color is direct and deterministic based on the provided conditions: "Maternal smoking status has a direct effect on brown eyes." Furthermore, the specific causal chain described states that "confounder active and smoking mother causes brown eyes."} \\
\\ 
\texttt{The question asks whether the infant would have **blue eyes** if the mother were a **nonsmoker** instead of a smoker, given that the unobserved confounder is active.} \\
\\ 
\texttt{1.  **Analyze the Nonsmoking Scenario**: If the mother is a nonsmoker, we must determine if the condition "confounder active and smoking mother" is met. Since the mother is not smoking, this specific causal trigger for brown eyes is broken.} \\
\texttt{2.  **Evaluate Alternative Causes**: The prompt explicitly states: "without any unmentioned factors or causal relationships." This is a critical constraint. There are no other variables mentioned (such as genetics, paternal factors, or environmental factors) that could cause blue eyes.} \\
\texttt{3.  **Logical Deduction**: In a closed system where "Smoking + Active Confounder" is the *only* defined cause for "Brown Eyes," and there are no other defined causes for "Blue Eyes," the absence of the cause for brown eyes does not automatically guarantee the presence of blue eyes unless a default state or alternative cause is specified. However, the question frames the outcome as a binary choice between the observed effect (brown eyes) and the alternative (blue eyes).} \\
\texttt{4.  **Causal Inference Principle**: In causal diagrams, if a variable (Brown Eyes) is caused by a specific set of parents (Smoking + Confounder), and we intervene to remove one parent (change to Nonsmoking), the effect (Brown Eyes) is blocked. Without an alternative mechanism defined in the world to produce Blue Eyes, the outcome is undefined or, in the context of a binary "Yes/No" regarding the specific alternative mentioned, the premise that "Blue Eyes" will occur is unsupported.} \\
\\ 
\texttt{However, looking closer at the phrasing "Would the infant has blue eyes...": This implies a counterfactual where the *only* other possibility is blue eyes. If the world contains *only* the factors listed, and "Smoking + Confounder" causes Brown Eyes, then removing the smoking factor removes the cause of Brown Eyes. In a strictly closed logical universe with no other causes for eye color mentioned, the infant's eye color would not be determined by the listed factors. But typically in such logic puzzles, if the specific condition for "Brown Eyes" is removed and no other condition for "Blue Eyes" is active, the answer to "Would it be Blue?" is **No** (because the condition for Blue Eyes was never established as active or default; the system simply lacks the cause for Brown Eyes, leaving the state undefined or null} \\
            }
        }
    }
    & \\ \\

    \theutterance \stepcounter{utterance}  
    & & \multicolumn{4}{p{0.6\linewidth}}{
        \cellcolor[rgb]{0.95,0.95,0.95}{
            \makecell[{{p{\linewidth}}}]{
                \texttt{\tiny{[P1$\leftrightarrow$P]}}
                \texttt{{[}SELF{-}FEEDBACK{]}} \\
\texttt{FORMAT: problem} \\
\texttt{RULES: problem} \\
\texttt{STRATEGY: problem} \\
\texttt{REASON: The answer fails to follow the required "Yes" or "No" prefix, contains excessive reasoning text, and provides a logically flawed conclusion regarding the closed{-}world assumption.} \\
            }
        }
    }
    & \\ \\

    \theutterance \stepcounter{utterance}  
    & \multicolumn{4}{p{0.6\linewidth}}{
        \cellcolor[rgb]{0.561,0.737,0.561}{
            \makecell[{{p{\linewidth}}}]{
                \texttt{\tiny{[P1$\rangle$GM]}}
                \texttt{No} \\
\\ 
\texttt{**Reasoning:**} \\
\\ 
\texttt{In this hypothetical world, the causal pathway from maternal smoking status to the infant's eye color is direct and deterministic based on the provided conditions: "Maternal smoking status has a direct effect on brown eyes." Furthermore, the specific causal chain described states that "confounder active and smoking mother causes brown eyes."} \\
\\ 
\texttt{The question asks whether the infant would have **blue eyes** if the mother were a **nonsmoker** instead of a smoker, given that the unobserved confounder is active.} \\
\\ 
\texttt{1.  **Analyze the Nonsmoking Scenario**: If the mother is a nonsmoker, the condition "confounder active and smoking mother" is no longer met. Consequently, the specific causal trigger for brown eyes is broken.} \\
\texttt{2.  **Evaluate Alternative Causes**: The prompt explicitly states: "without any unmentioned factors or causal relationships." This is a critical constraint. There are no other variables mentioned (such as genetics, paternal factors, or environmental factors) that could cause blue eyes.} \\
\texttt{3.  **Logical Deduction**: In a closed system where "Smoking + Active Confounder" is the *only* defined cause for "Brown Eyes," and there are no other defined causes for "Blue Eyes," the absence of the cause for brown eyes does not automatically guarantee the presence of blue eyes. However, the question frames the outcome as a binary choice between the observed effect (brown eyes) and the alternative (blue eyes). If the specific condition for "Brown Eyes" is removed and no other condition for "Blue Eyes" is active, the premise that "Blue Eyes" will occur is unsupported. Therefore, the answer to "Would the infant have blue eyes?" is **No**, as the necessary conditions for blue eyes were never established in this closed world.} \\
            }
        }
    }
    & & \\ \\

    \theutterance \stepcounter{utterance}  
    & & & \multicolumn{2}{p{0.3\linewidth}}{
        \cellcolor[rgb]{0.9,0.9,0.9}{
            \makecell[{{p{\linewidth}}}]{
                \texttt{\tiny{[GM$|$GM]}}
                \texttt{No} \\
            }
        }
    }
    & & \\ \\

    \theutterance \stepcounter{utterance}  
    & & & \multicolumn{2}{p{0.3\linewidth}}{
        \cellcolor[rgb]{0.9,0.9,0.9}{
            \makecell[{{p{\linewidth}}}]{
                \texttt{\tiny{[GM$|$GM]}}
                \texttt{no} \\
            }
        }
    }
    & & \\ \\

    \theutterance \stepcounter{utterance}  
    & & & \multicolumn{2}{p{0.3\linewidth}}{
        \cellcolor[rgb]{0.9,0.9,0.9}{
            \makecell[{{p{\linewidth}}}]{
                \texttt{\tiny{[GM$|$GM]}}
                \texttt{game\_result = WIN} \\
            }
        }
    }
    & & \\ \\

\end{supertabular}
}

\end{document}
